\section{生成式视觉语言模型设计}

语义分类实验已经表明 EEG 能够辅助视觉特征提升识别效果，接下来我们把 EEG-enhanced vision tokens 接入生成式视觉语言模型，搭建一条从退化图像和 EEG 到自然语言描述的完整生成通路。

\subsection{生成式描述任务的设计目标}

我们这个生成式描述任务的目标很明确：给定一张退化图像以及与之配对的一条 EEG trial，模型要输出一段自然语言的 caption 来描述图像内容。和第 5 章的语义分类任务不一样，这里模型不输出固定的类别标签，而是要生成自由形式的描述文本。

需要澄清一点，这个任务的目标并不是"从脑电信号里直接读出任意语言"。在我们大约 40 个类别、12000 条 trial 的数据规模下，设计思路是把 EEG 当作一种视觉语义增强信号来用，利用 EEG 携带的额外神经表征信息，帮助生成式视觉语言模型在图像退化条件下输出比纯视觉输入更准确的描述。

\subsection{轻量级 GRU 图像描述基线}

在接入大模型之前，我们先搭了一个轻量的 GRU-based autoregressive caption decoder 作为生成基线。这个 decoder 拿 VTF 输出的 enhanced vision tokens 做平均池化，把池化向量作为初始隐状态，然后用单层 GRU 逐步解码出文本 token。模型参数量大概 2.8M，训练数据用的是 BLIP 生成的伪 caption。

GRU baseline 能稳定地生成 free-form caption，好处是训练快、不依赖大型预训练模型、输出始终是长度自然的文本。不过 GRU decoder 缺少丰富的语言先验知识，生成的描述在语法自然度和概念准确度上明显比预训练视觉语言模型差。GRU baseline 的主要意义在于验证了 EEG-enhanced vision tokens 确实能驱动自回归 caption 生成。

\subsection{预训练生成模型方案比较}

为了确定最优的生成方案，我们分别试了五种技术路线，汇总在表~\ref{tab:route-comparison} 里。实验阶段为了方便记录，用代号来标记各个方案。

\begin{table}[htbp]
  \centering
  \caption{生成式视觉语言模型方案对比}
  \label{tab:route-comparison}
  \begin{tabularx}{\linewidth}{l l X}
    \toprule
    方案代号 & 输入与桥接方式 & 主要发现 \\
    \midrule
    方案一 & enhanced tokens 池化后投影为 inputs\_embeds，输入 Qwen2.5 + LoRA & 纯文本 LLM 缺乏视觉-语言对齐能力，有效输出率低 \\
    方案二 & Q-Former/Perceiver bridge 桥接视觉与语言 & 在小样本条件下效果有限，需要更大的训练数据 \\
    方案三 & LLaVA-style MLP projector & 与 Qwen2-VL adapter 的工程兼容性不佳 \\
    方案四 & BLIP-2 style Q-Former & 清洁图像上表现尚可，退化条件下不稳定 \\
    方案五 & enhanced tokens + Qwen2-VL adapter/prefix + LoRA & 最终采用的生成方案 \\
    \bottomrule
  \end{tabularx}
\end{table}

综合对比下来，方案五（Qwen2-VL adapter）在有效输出率、caption class-hit 和 real-control gap 这三个指标上都明显好于前四种方案。Qwen2-VL 本身是原生视觉语言模型，它的视觉编码模块在预训练阶段就已经和语言模型做过充分的图文对齐，所以在小样本微调的场景下能更好地利用 EEG-enhanced visual tokens。

\subsection{基于 Qwen2-VL 的最终生成方案}

我们最终采用的生成方案结构如图~\ref{fig:route5} 所示。

\placeholderfigure{基于 Qwen2-VL 的生成式视觉语言模型架构图}{route5-qwenvl-placeholder}{5.0cm}

核心流程是这样的：(1) 退化图像经过 CLIP ViT 得到 visual tokens，和 EEG tokens 一起通过 VTF cross-attention 融合得到增强视觉 token；(2) 增强视觉 token 先经过我们自己实现的轻量 prefix adapter，映射到 Qwen2-VL 能接收的嵌入空间，然后和语义提示文本一起输入生成式解码器；(3) 输入时额外附加语义提示信息（A2 预测的最可能类别）和退化类型标签，以文本 token 形式插入到生成 prompt 中；(4) Qwen2-VL-2B-Instruct 解码器在 teacher-forcing 模式下做自回归训练，测试时以 ``corruption\_type:'' 作为前缀来触发生成。

\subsection{LoRA 微调与训练配置}

生成式模型用 LoRA 对 Qwen2-VL 的 Query 和 Value 投影矩阵做低秩适配，我们设的秩 $r=8$，缩放因子 $\alpha=16$。可训练参数大约 1.3M，跟模型总参数量比很小。训练的时候冻结了所有原始模型权重。另外我们还对比了 $r=16$ 的配置。

训练目标方面，我们试了两种策略：(1) 只用类别名称作为目标——直接以类别正确性为优化目标；(2) 类别名称配合过滤后的 BLIP 伪 caption 作目标——生成的文本更自然，但 BLIP 伪 caption 里可能带噪声，会影响类别准确性。实验表明，类别目标更有利于提高 caption 类别命中率，而类别配合 BLIP 伪 caption 目标更有利于提高文本自然度和有效输出率。我们分别报告了两种目标的效果，最终采用 Qwen2-VL LoRA $r=8$ 配合多候选重排序作为主要的生成结果。

\subsection{多候选重排序策略}

推理的时候，我们对每个输入生成 $N=5$ 个候选 caption（通过调整温度参数、top-p 等解码策略实现），先根据候选 caption 的有效性（不能是空、不能是 URL 或代码）、重复率和长度做一轮初步过滤，然后算每个候选跟 A2/VTF 预测的 top-k 语义概念或 CLIP 文本原型之间的相似度，挑出语义一致性最高的那个作为最终输出。整个过程不依赖测试集的真实类别标签，caption 类别命中率只在最终评估时才用到。多候选重排序不改模型参数，但明显提升了有效输出率，定性展示效果也更好了。
